% GENERATED FILE. Edit canonical lessons, then run npm run generate:books.
% canonical-source-hash: fnv1a64:ba44eb993659492c
% canonical-lessons: ML-C502-tottam, ML-C502-paccappu, ML-C502-talirkkuka, ML-C502-koliyuka, ML-C502-vettimarruka

\chapter{Plantations and Trees}
\label{ch:ml-plantations-and-trees}
\hlchaptermodality{502}

\begin{chapteropening}
\textbf{By the end of this chapter:} \emph{I can say a plantation, greenery, to sprout, to fall off and to cut down.}

Complete the last lesson of chapter 502: I can say a plantation, greenery, to sprout, to fall off and to cut down.
\end{chapteropening}

\section[\texorpdfstring{\ml{തോട്ടം} (tōṭṭaṁ)}{തോട്ടം (tōṭṭaṁ)}]{\ml{തോട്ടം} (tōṭṭaṁ) — a plantation; a garden}

\label{lesson:ML-C502-tottam}



\begin{quote}
Before the new one: say the Malayalam for an entry fee, then the Malayalam for to drift in.
\end{quote}

\subsection*{You'll want to know: \ml{തോട്ടം}}
\textbf{\ml{തോട്ടം}} — \emph{tōṭṭaṁ} — \textquotedblleft{}a plantation; a garden\textquotedblright{}.

\begin{quote}
\ml{മലയിൽ} \ml{വലിയ} \ml{തേയിലത്തോട്ടങ്ങൾ} \ml{ഉണ്ട്}.
\end{quote}

\begin{quote}
\emph{malayil valiya tēyilattōṭṭaṅṅaḷ uṇṭŭ.}
\end{quote}

\begin{quote}
There are big tea plantations in the hills.
\end{quote}

Plural: \ml{തോട്ടങ്ങൾ}, \emph{tōṭṭaṅṅaḷ}.

You met it in \ml{പൂന്തോട്ടം}, a flower garden.

\begin{cousinweb}
\textbf{\ml{തോട്ടം}} (\emph{tōṭṭaṁ}) is native Dravidian, \textquotedblleft{}a garden\textquotedblright{}, from a root that means to dig (DEDR 3549).
\end{cousinweb}

\begin{grammarlens}[title={What you've built}]
One more word for getting around.
\end{grammarlens}

\subsection*{Your turn}
\begin{itemize}
\raggedright
  \item \emph{Say it:} \emph{tōṭṭaṁ}
  \item \emph{Say it:} \emph{tōṭṭaṁ}, once more
  \item \emph{Say it:} \emph{malayil valiya tēyilattōṭṭaṅṅaḷ uṇṭŭ.} — There are big tea plantations in the hills.
  \item \emph{Recall:} say the Malayalam for an entry fee, then the Malayalam for to drift in, then say \emph{tōṭṭaṁ} again
\end{itemize}

\begin{tcolorbox}[breakable,colback=teal!4,colframe=teal!35!black,title={Before you move on}]
What does \ml{തോട്ടം} mean? (\textquotedblleft{}A plantation; a garden\textquotedblright{}.) Say it once more.
\end{tcolorbox}
\section[\texorpdfstring{\ml{പച്ചപ്പ്} (paccappŭ)}{പച്ചപ്പ് (paccappŭ)}]{\ml{പച്ചപ്പ്} (paccappŭ) — greenery}

\label{lesson:ML-C502-paccappu}



\begin{quote}
Before the new one: say the Malayalam for to drift in, then the Malayalam for a plantation.
\end{quote}

\subsection*{You'll want to know: \ml{പച്ചപ്പ്}}
\textbf{\ml{പച്ചപ്പ്}} — \emph{paccappŭ} — \textquotedblleft{}greenery\textquotedblright{}.

\begin{quote}
\ml{മഴക്കാലത്ത്} \ml{എവിടെയും} \ml{പച്ചപ്പാണ്}.
\end{quote}

\begin{quote}
\emph{maḻakkālattŭ eviṭeyuṁ paccappāṇŭ.}
\end{quote}

\begin{quote}
In the rainy season there is greenery everywhere.
\end{quote}

The noun of \ml{പച്ച}, green, with -\ml{പ്പ്}.

\begin{cousinweb}
\textbf{\ml{പച്ചപ്പ്}} (\emph{paccappŭ}) is native Dravidian, the noun of \ml{പച്ച}, green (DEDR 3821).
\end{cousinweb}

\begin{grammarlens}[title={What you've built}]
One more word for getting around.
\end{grammarlens}

\subsection*{Your turn}
\begin{itemize}
\raggedright
  \item \emph{Say it:} \emph{paccappŭ}
  \item \emph{Say it:} \emph{paccappŭ}, once more
  \item \emph{Say it:} \emph{maḻakkālattŭ eviṭeyuṁ paccappāṇŭ.} — In the rainy season there is greenery everywhere.
  \item \emph{Recall:} say the Malayalam for to drift in, then the Malayalam for a plantation, then say \emph{paccappŭ} again
\end{itemize}

\begin{tcolorbox}[breakable,colback=teal!4,colframe=teal!35!black,title={Before you move on}]
What does \ml{പച്ചപ്പ്} mean? (\textquotedblleft{}Greenery\textquotedblright{}.) Say it once more.
\end{tcolorbox}
\section[\texorpdfstring{\ml{തളിർക്കുക} (taḷirkkuka)}{തളിർക്കുക (taḷirkkuka)}]{\ml{തളിർക്കുക} (taḷirkkuka) — to sprout; to put out leaves}

\label{lesson:ML-C502-talirkkuka}



\begin{quote}
Before the new one: say the Malayalam for a plantation, then the Malayalam for greenery.
\end{quote}

\subsection*{You'll want to know: \ml{തളിർക്കുക}}
\textbf{\ml{തളിർക്കുക}} — \emph{taḷirkkuka} — \textquotedblleft{}to sprout; to put out leaves\textquotedblright{}.

\begin{quote}
\ml{മഴ} \ml{പെയ്തപ്പോൾ} \ml{ഉണങ്ങിയ} \ml{ചെടികൾ} \ml{തളിർത്തു}.
\end{quote}

\begin{quote}
\emph{maḻa peytappōḷ uṇaṅṅiya ceṭikaḷ taḷirttu.}
\end{quote}

\begin{quote}
When the rain fell, the dry plants sprouted.
\end{quote}

Past: \ml{തളിർത്തു}, \emph{taḷirttu}.

\ml{തളിര്} is a tender new leaf.

\begin{cousinweb}
\textbf{\ml{തളിർക്കുക}} (\emph{taḷirkkuka}) is native Dravidian, \textquotedblleft{}to bud, to sprout\textquotedblright{}, from \ml{തളിര്}, a new leaf (DEDR 3131).
\end{cousinweb}

\begin{grammarlens}[title={What you've built}]
One more word for getting around.
\end{grammarlens}

\subsection*{Your turn}
\begin{itemize}
\raggedright
  \item \emph{Say it:} \emph{taḷirkkuka}
  \item \emph{Say it:} \emph{taḷirkkuka}, once more
  \item \emph{Say it:} \emph{maḻa peytappōḷ uṇaṅṅiya ceṭikaḷ taḷirttu.} — When the rain fell, the dry plants sprouted.
  \item \emph{Recall:} say the Malayalam for a plantation, then the Malayalam for greenery, then say \emph{taḷirkkuka} again
\end{itemize}

\begin{tcolorbox}[breakable,colback=teal!4,colframe=teal!35!black,title={Before you move on}]
What does \ml{തളിർക്കുക} mean? (\textquotedblleft{}To sprout; to put out leaves\textquotedblright{}.) Say it once more.
\end{tcolorbox}
\section[\texorpdfstring{\ml{കൊഴിയുക} (koḻiyuka)}{കൊഴിയുക (koḻiyuka)}]{\ml{കൊഴിയുക} (koḻiyuka) — to fall off; to drop}

\label{lesson:ML-C502-koliyuka}



\begin{quote}
Before the new one: say the Malayalam for greenery, then the Malayalam for to sprout.
\end{quote}

\subsection*{You'll want to know: \ml{കൊഴിയുക}}
\textbf{\ml{കൊഴിയുക}} — \emph{koḻiyuka} — \textquotedblleft{}to fall off; to drop\textquotedblright{}.

\begin{quote}
\ml{വേനലിൽ} \ml{മരങ്ങളുടെ} \ml{ഇലകൾ} \ml{കൊഴിഞ്ഞു}.
\end{quote}

\begin{quote}
\emph{vēnalil maraṅṅaḷuṭe ilakaḷ koḻiññu.}
\end{quote}

\begin{quote}
In the dry season the leaves fell from the trees.
\end{quote}

Past: \ml{കൊഴിഞ്ഞു}, \emph{koḻiññu}.

For leaves, fruit, petals and hair; a person falls with \ml{വീഴുക}.

\begin{cousinweb}
\textbf{\ml{കൊഴിയുക}} (\emph{koḻiyuka}) is native Dravidian, \textquotedblleft{}to fall out, to drop, as fruit, leaves or hair\textquotedblright{} (DEDR 1825).
\end{cousinweb}

\begin{grammarlens}[title={What you've built}]
One more word for getting around.
\end{grammarlens}

\subsection*{Your turn}
\begin{itemize}
\raggedright
  \item \emph{Say it:} \emph{koḻiyuka}
  \item \emph{Say it:} \emph{koḻiyuka}, once more
  \item \emph{Say it:} \emph{vēnalil maraṅṅaḷuṭe ilakaḷ koḻiññu.} — In the dry season the leaves fell from the trees.
  \item \emph{Recall:} say the Malayalam for greenery, then the Malayalam for to sprout, then say \emph{koḻiyuka} again
\end{itemize}

\begin{tcolorbox}[breakable,colback=teal!4,colframe=teal!35!black,title={Before you move on}]
What does \ml{കൊഴിയുക} mean? (\textquotedblleft{}To fall off; to drop\textquotedblright{}.) Say it once more.
\end{tcolorbox}
\section[\texorpdfstring{\ml{വെട്ടിമാറ്റുക} (veṭṭimāṟṟuka)}{വെട്ടിമാറ്റുക (veṭṭimāṟṟuka)}]{\ml{വെട്ടിമാറ്റുക} (veṭṭimāṟṟuka) — to cut down; to clear away}

\label{lesson:ML-C502-vettimarruka}



\begin{quote}
Before the new one: say the Malayalam for to sprout, then the Malayalam for to fall off.
\end{quote}

\subsection*{You'll want to know: \ml{വെട്ടിമാറ്റുക}}
\textbf{\ml{വെട്ടിമാറ്റുക}} — \emph{veṭṭimāṟṟuka} — \textquotedblleft{}to cut down; to clear away\textquotedblright{}.

\begin{quote}
\ml{വഴി} \ml{വീതി} \ml{കൂട്ടാൻ} \ml{അവർ} \ml{മരങ്ങൾ} \ml{വെട്ടിമാറ്റി}.
\end{quote}

\begin{quote}
\emph{vaḻi vīti kūṭṭān avar maraṅṅaḷ veṭṭimāṟṟi.}
\end{quote}

\begin{quote}
They cut down the trees to widen the road.
\end{quote}

Past: \ml{വെട്ടിമാറ്റി}, \emph{veṭṭimāṟṟi}.

\ml{വെട്ടി}, having cut, plus \ml{മാറ്റുക}, to remove.

\begin{cousinweb}
\textbf{\ml{വെട്ടിമാറ്റുക}} (\emph{veṭṭimāṟṟuka}) is a native Dravidian compound of \ml{വെട്ടുക}, to cut with an axe (DEDR 5478), and \ml{മാറ്റുക}, to move away (DEDR 4834).
\end{cousinweb}

\begin{grammarlens}[title={What you've built}]
One more word for getting around.
\end{grammarlens}

\subsection*{Your turn}
\begin{itemize}
\raggedright
  \item \emph{Say it:} \emph{veṭṭimāṟṟuka}
  \item \emph{Say it:} \emph{veṭṭimāṟṟuka}, once more
  \item \emph{Say it:} \emph{vaḻi vīti kūṭṭān avar maraṅṅaḷ veṭṭimāṟṟi.} — They cut down the trees to widen the road.
  \item \emph{Recall:} say the Malayalam for to sprout, then the Malayalam for to fall off, then say \emph{veṭṭimāṟṟuka} again
\end{itemize}

\begin{tcolorbox}[breakable,colback=teal!4,colframe=teal!35!black,title={Before you move on}]
What does \ml{വെട്ടിമാറ്റുക} mean? (\textquotedblleft{}To cut down; to clear away\textquotedblright{}.) Say it once more.
\end{tcolorbox}
